% Source: Weinberg GR, physical PDF pp. 67--69
%         (printed pp. 41--43).
% Coverage: Section 2.7, Electrodynamics; equations
%           (2.7.1)--(2.7.13).
% Figures/tables/footnotes: none.
% Status: source-reviewed and compile-clean.
% Uncertainties: none.

\subsection{Electrodynamics}\label{sec:2.7}

Maxwell's equations for the electric and magnetic fields \(\mathbf E,\mathbf B\)
produced by a given charge density \(\rho\) and current density \(\mathbf J\)
are
\begin{equation}
  \boldsymbol{\nabla}\cdot\mathbf E=\rho ,
  \tag{2.7.1}\label{eq:2.7.1}
\end{equation}
\begin{equation}
  \boldsymbol{\nabla}\times\mathbf B
  =\frac{\partial\mathbf E}{\partial t}+\mathbf J ,
  \tag{2.7.2}\label{eq:2.7.2}
\end{equation}
\begin{equation}
  \boldsymbol{\nabla}\cdot\mathbf B=0 ,
  \tag{2.7.3}\label{eq:2.7.3}
\end{equation}
\begin{equation}
  \boldsymbol{\nabla}\times\mathbf E
  =-\frac{\partial\mathbf B}{\partial t}.
  \tag{2.7.4}\label{eq:2.7.4}
\end{equation}
To uncover the Lorentz transformation properties of \(\mathbf E\) and
\(\mathbf B\), we introduce a matrix \(F^{\mu\nu}\), defined by
\begin{equation}
  F^{0i}=E_i,\qquad
  F^{ij}=\varepsilon^{ijk}B_k,\qquad
  F^{\mu\nu}=-F^{\nu\mu},
  \tag{2.7.5}\label{eq:2.7.5}
\end{equation}
with \(\varepsilon^{123}=+1\). Then Eqs.~\eqref{eq:2.7.1} and
\eqref{eq:2.7.2} can be written as
\begin{equation}
  \partial_\mu F^{\mu\nu}=-J^\nu
  \tag{2.7.6}\label{eq:2.7.6}
\end{equation}
(recall that \(J^0=\rho\)), whereas Eqs.~\eqref{eq:2.7.3} and
\eqref{eq:2.7.4} give
\begin{equation}
  \varepsilon^{\mu\nu\rho\sigma}\partial_\nu F_{\rho\sigma}=0 ,
  \tag{2.7.7}\label{eq:2.7.7}
\end{equation}
where \(\varepsilon^{\mu\nu\rho\sigma}\) is the Levi-Civita tensor defined in
Section~\ref{sec:2.5}, and \(F_{\rho\sigma}\) is the covariant tensor defined
as usual by
\[
  F_{\rho\sigma}
  =\eta_{\rho\mu}\eta_{\sigma\nu}F^{\mu\nu}.
\]
Since \(J^\mu\) is a four-vector, we conclude that \(F^{\mu\nu}\) is a tensor,
\begin{equation}
  F'^{\mu\nu}
  =\Lambda^\mu{}_\rho\Lambda^\nu{}_\sigma F^{\rho\sigma},
  \tag{2.7.8}\label{eq:2.7.8}
\end{equation}
because if \(F^{\mu\nu}\) is a solution of Eqs.~\eqref{eq:2.7.6} and
\eqref{eq:2.7.7}, then Eq.~\eqref{eq:2.7.8} will be a solution in a
Lorentz-transformed coordinate system.

The electromagnetic force on a charged particle is
\begin{equation}
  f^\mu=eF^\mu{}_\nu\frac{dx^\nu}{d\tau}.
  \tag{2.7.9}\label{eq:2.7.9}
\end{equation}
That this is correct may be seen by repeating the arguments of
Section~\ref{sec:2.3}. Equation \eqref{eq:2.7.9} is correct in a reference
system in which the particle is at rest, because in this frame it gives
\(\mathbf f=e\mathbf E,\ f^0=0\), and it transforms like a four-vector, so it
is correct for all velocities. Note incidentally that Eqs.~\eqref{eq:2.7.9}
and \eqref{eq:2.4.2} give
\[
  \frac{d\mathbf p}{dt}
  =e(\mathbf E+\mathbf v\times\mathbf B),
\]
so the formula for magnetic force follows as a consequence of special
relativity.

There is a useful alternate form to the homogeneous equations
\eqref{eq:2.7.7}:
\begin{equation}
  \partial_\mu F_{\nu\rho}
  +\partial_\nu F_{\rho\mu}
  +\partial_\rho F_{\mu\nu}=0 .
  \tag{2.7.10}\label{eq:2.7.10}
\end{equation}
Note that for \(\mu,\nu,\rho\) all different, Eq.~\eqref{eq:2.7.10} is the
same as Eq.~\eqref{eq:2.7.7}; for instance, setting the free index to \(0\) in
Eq.~\eqref{eq:2.7.7} gives the same result as setting
\(\mu\nu\rho=123\) in Eq.~\eqref{eq:2.7.10}. On the other hand, for two
indices equal, Eq.~\eqref{eq:2.7.10} is an identity; for instance, if
\(\nu=\rho\), then Eq.~\eqref{eq:2.7.10} reads
\[
  \partial_\mu F_{\nu\nu}
  +\partial_\nu F_{\nu\mu}
  +\partial_\nu F_{\mu\nu}=0
  \qquad(\text{not summed}),
\]
and this is identically true because \(F_{\mu\nu}=-F_{\nu\mu}\).

Equation \eqref{eq:2.7.7} allows us to represent \(F_{\mu\nu}\) as a
``curl'' of a four-vector \(A_\nu\):
\begin{equation}
  F_{\mu\nu}=\partial_\mu A_\nu-\partial_\nu A_\mu .
  \tag{2.7.11}\label{eq:2.7.11}
\end{equation}
(See Section~4.11.) We can change \(A_\mu\) by a term
\(\partial_\mu\phi\) without affecting \(F_{\mu\nu}\), so \(A_\mu\) may be
defined so that
\begin{equation}
  \partial^\mu A_\mu=0 .
  \tag{2.7.12}\label{eq:2.7.12}
\end{equation}
With Eqs.~\eqref{eq:2.7.11} and \eqref{eq:2.7.12}, the rest of Maxwell's
equations reduce to
\begin{equation}
  \Box A_\mu=-J_\mu .
  \tag{2.7.13}\label{eq:2.7.13}
\end{equation}
